% =====================================================================
%  Glosario de términos
% =====================================================================
\chapter*{GLOSARIO DE TÉRMINOS}
\addcontentsline{toc}{chapter}{GLOSARIO DE TÉRMINOS}

\begin{description}
  \item[Ahorro de selecciones.] Porcentaje de selecciones que se evitan al escribir un texto con
  ayuda de la predicción de palabras, respecto a escribirlo carácter por carácter.

  \item[Artefacto.] Componente del registro de EEG que no proviene de la actividad cerebral, como
  los parpadeos, la actividad muscular o la interferencia de la red eléctrica.

  \item[AUC (área bajo la curva ROC).] Métrica que indica la probabilidad de que un clasificador
  asigne un puntaje mayor a una época con P300 que a una sin ella. Un valor de 0.5 corresponde a un
  clasificador aleatorio y un valor de 1 a una separación perfecta.

  \item[BCI (interfaz cerebro-computadora).] Sistema de comunicación en el que los mensajes o
  comandos del usuario se obtienen de su actividad cerebral, sin pasar por los nervios periféricos
  ni por los músculos.

  \item[Calibración.] Sesión en la que el usuario fija su atención en caracteres indicados por el
  sistema mientras se registran sus señales, con las que se entrena su propio modelo de detección.

  \item[CNN (red neuronal convolucional).] Red neuronal que aplica filtros de pesos aprendibles
  sobre la entrada para extraer características, adecuada para datos como imágenes o series de
  tiempo.

  \item[EEG (electroencefalograma).] Registro de la actividad eléctrica del cerebro tomado por medio
  de electrodos colocados en la superficie de la cabeza.

  \item[EEGNet.] Red neuronal convolucional compacta diseñada para clasificar señales de EEG en
  distintos paradigmas de BCI, entre ellos la P300.

  \item[ELA (esclerosis lateral amiotrófica).] Enfermedad neurológica que en sus etapas avanzadas
  puede provocar la pérdida casi total del control muscular voluntario, incluido el del habla.

  \item[Época.] Segmento de la señal de EEG que inicia en el instante de una intensificación y que
  constituye la muestra de entrada al clasificador.

  \item[ERP (potencial relacionado con eventos).] Del inglés \textit{Event-Related Potential}.
  Variación del voltaje del EEG que aparece ligada en el tiempo a un evento específico, como un
  estímulo sensorial, una respuesta motora o una operación cognitiva. Se registra como una serie de
  ondas positivas y negativas llamadas componentes, que se nombran con la letra P o N según su
  polaridad, seguida de su latencia aproximada en milisegundos. Debido a que su amplitud es mucho
  menor que la de la actividad de fondo del EEG, para observarlo se suelen promediar varias
  repeticiones del mismo estímulo. La P300 es la componente de este tipo que aprovecha el presente
  trabajo.

  \item[Exactitud por carácter.] Proporción de caracteres que el sistema selecciona correctamente
  respecto al total de caracteres que el usuario intentó escribir.

  \item[Hiperparámetro.] Valor que define la configuración de un modelo o de su entrenamiento, como
  el número de filtros o la tasa de aprendizaje, y que no se aprende a partir de los datos.

  \item[Intensificación.] Iluminación momentánea de una fila o de una columna de la matriz del P300
  Speller.

  \item[ISI (intervalo entre estímulos).] Tiempo que la matriz permanece sin iluminar entre el fin
  de una intensificación y el inicio de la siguiente.

  \item[ITR (tasa de transferencia de información).] Cantidad de información que el usuario
  transmite por unidad de tiempo, expresada en bits por minuto, que combina la exactitud y la
  velocidad de las selecciones.

  \item[Modelo de lenguaje.] Modelo que asigna una probabilidad a cada secuencia de palabras, lo que
  permite estimar cuáles son las palabras más probables dado lo que se ha escrito.

  \item[P300.] Componente positiva de los ERP que aparece alrededor de 300~ms después de un estímulo
  poco frecuente y relevante para la persona.

  \item[P300 Speller.] Sistema de deletreo en el que las filas y columnas de una matriz de
  caracteres se iluminan en orden aleatorio, y el carácter que el usuario desea se identifica a
  partir de la P300 que provoca su fila y su columna.

  \item[Paradigma del evento raro (\textit{oddball}).] Procedimiento en el que se presenta una
  secuencia de estímulos de dos categorías, una de ellas poco frecuente, y se pide a la persona
  atender solo a los estímulos de esa categoría, que son los que provocan la P300.

  \item[Repetición.] Secuencia en la que cada fila y cada columna de la matriz se intensifican una
  vez.

  \item[Sistema internacional 10-20.] Estándar que define la posición de los electrodos de EEG sobre
  la cabeza a partir de distancias del 10\% y del 20\% entre puntos de referencia anatómicos.

  \item[Sobreajuste.] Situación en la que un modelo memoriza los datos de entrenamiento y pierde la
  capacidad de desempeñarse correctamente con datos nuevos.
\end{description}
